\section{Week 4: Incarnation of the Logos}

\begin{quotationx}
The tendency is certainly accentuated, if not prevalent, amongst contemporary Hermeticists to concern themselves more
with the “Cosmic Christ” or the “Logos” than with the human person of the “Son of Man”, Jesus of Nazareth. More
importance is attributed to the divine and abstract aspect of the God-Man than to his human and concrete aspect.

It was contact with the person of Jesus Christ which opened up the current of miracles and conversions. And it is the
same even today. \flright{\textsc{Valentin Tomberg}, \emph{Letter VIII: Justice}}

\end{quotationx}

With these words, Tomberg is warning us not to forget about the first coming of Jesus in the flesh, regarded as somehow
inferior to an esoteric interpretation. A fortiori, the Hermetist's goal is not to create an alternative or
“better” religion. Nevertheless, there is always a stream of “New Age” gurus who claim something similar. For example,
one such popular guru claims to have discovered the real meaning of all the religions, viz., what the Buddha “really”
taught or what Christ “really” taught. He then claims that the religions have distorted those teachings and offer no
authentic path. Although he came to that realization spontaneously, he will teach you certain “modalities” for a hefty
price to reach the same realization. This is the sin of simony, the notion that spiritual enlightenment is a commodity
that can be bought and sold.

The idea of the Logos was not unknown to pagan philosophers and Hermetists prior to the first Christmas. However, it is
the fact of the Incarnation that matters most, as St John pointed out in the remarkable claim that the Logos became
flesh. So Jesus is not only the fulfillment of the Mosaic law, He is also the fulfillment of the natural law. This is
made clear by the visit of the Magi.

Tomberg makes us wrestle with a philosophical conundrum. The thinking mind, restricted to dianoia, knows essences, and
the Logos is “the fundamental universal [or essence] of the world”. And Jesus Christ is then the “particular of
particulars”. Some minds, like that of the new age guru, see that a representing a limitation on their thought; hence
they resort to a sort of Docetism which denies the need for the physical, including a birth, visible church, sacraments
and so on. It is a small step, then, to reach the conclusion that there is no need for the purification of the head and
the heart in order to reach higher states.

Since for God, essence and existence are One, to know God is to know both his essence and existence. Tomberg explains
that

\begin{quotex}
Christian Hermeticism itself can only be knowledge of the universal which is revealed in the particular. 

\end{quotex}

Hence, the Christian Hermetist “aspires to mystical experience of the communion of beings through love”. Thus he seeks
spiritual friendships in the particular.

Yet, not unlike the pagan Hermetists — his precursors — or even the new ager perhaps, he also
seeks the mystical experience of communion with the Logos, i.e., the knowledge of the universal.

\paragraph{Spiritual Beings}
Fr. Reginald Garrigou-Lagrange tells us that the angels know intuitively, not rationally. Each higher level of angel
understands more through the knowledge of ever more encompassing principles. Tomberg asserts:

\begin{quotex}
For Hermeticism there are no “principles”, “laws”, and “ideas” which exist outside of individual beings, not as
structural traits of their nature, but as entities separated and independent from it. 

\end{quotex}

This makes perfect sense, since knowing and being are one. If an angel, then, “knows” a certain principle, it is ipso
fact the embodiment of that principle. Ideas have no power on their own, they are purely passive. An idea has effects
only when it is immanent in a being. We can choose to understand our environment as an abstraction, the mere interplay
of impersonal forces. Or else, we can choose to understand it as a great drama of personal forces.

A recent episode of the Vikings series on the History Channel had an interesting scene. Rollo was a Viking warrior who
converted to Catholicism and was rewarded with the Duchy of Normandy. Unable to totally forget his pagan past, he
explained to his wife, “When you hear thunder, it is only thunder. But when I hear thunder, I hear the sound of Thor.”
He was still at the level of original participation in Owen Barfield's sense. Can we recover that state of
original participation?

Tomberg tells us we must “love our pagan past”, so perhaps we can learn something from Rollo. Now this is not a new
teaching, but actually something we forgot. So perhaps we can try to remember. The mystic visionary, Catherine
Emmerich, saw that the world was populated with angels: each country, city, diocese, and parish has its own guardian
angel. Fr. Ripperger, in a youtube video\footnote{\url{https://youtu.be/mpJgVAs02Dc?si=DRtUFauKacdQ_HyN}}, reminds us furthermore that each generation has “generational spirits”, not
all benign, as a sort of Zeitgeist.

If we can overcome the Zeitgeist of scientism, we can meditate on our role in the cosmic hierarchy. See yourself in
relation to your family, parish or other community, nation, Church, then ascending through the angelic hierarchy. And
when you get to the Logos, see also the Baby in the manger.

\flright{\itshape Posted on 2022-12-16 by Cologero}
